\section{Integer memory and the failure of a naive partition formula}
\label{sec:integer-revelation}
The divisibility assumptions in Theorem~\ref{thm:joint-law} can be
removed, but a balanced-partition Bayes bound must not be mistaken for
a minimax equality when a persistent random partition is not free.
This distinction leads to another exact finite-resource transition.

Keep the revelation interface, but allow any integers $m\ge2$ and
$1\le K\le m$. Put $\rho=1-(1-\alpha)^N$, and denote the genuine
$K$-label minimax value by $V^{\rm er}_{N,m}(K)$. Write
$m=Kq+s$, $0\le s<K$, and set
\begin{equation}\label{eq:balanced-integer}
 B_{m,K}=1-\frac{(K-s)q^2+s(q+1)^2}{m^2}.
\end{equation}
Here and below all-erasure initialization randomness is a state
initialization, not an unrecorded seed used again at validation.

\begin{proposition}[Arbitrary integer alphabet bounds and low-reveal equality]
\label{prop:integer-erasure}
The optimal uniform-prior Bayes score is $\rho B_{m,K}$. The minimax
value satisfies
\begin{equation}\label{eq:integer-bounds}
 V^{\rm er}_{N,m}(K)\le
 \min\left\{\rho B_{m,K},\;1-\frac{\lceil m/K\rceil}{m}\right\}.
\end{equation}
If $\rho\le m/(m+s)$, equality with the first term holds. At perfect
revelation $\rho=1$ the exact value is the second term. In particular,
when $s>0$ the uniform-prior Bayes optimum at perfect revelation is
strictly larger than the minimax value.
\end{proposition}
\begin{proof}
The all-erasure decomposition \eqref{eq:erasure-channel} holds without
any divisibility condition. Its parameter-independent part has zero
uniform-prior score. Optimizing the other part over responses gives a
convex function of its stochastic channel. A deterministic channel with
fiber sizes $n_j$ has Bayes score $1-\sum_j(n_j/m)^2$. If two nonempty
fiber sizes differ by at least two, moving one index from the larger
to the smaller strictly decreases the sum of squares. Allowing empty
fibers and repeating gives sizes $q$ and $q+1$, proving the Bayes bound.
It is attained by this partition and its complementary events, with
any parameter-independent initialization. This proves the first upper
bound as well as the asserted Bayes optimum.

For perfect revelation, fix any $K$ response vectors. For each known
parameter choose a response maximizing its score; this deterministic
assignment cannot decrease any parameter's payoff, even when the
original encoder was randomized. Some assigned fiber has size at least
$\lceil m/K\rceil$. Averaging its parameters, the score of their common
response is at most the total variation between the uniform target
and the uniform law on that fiber, namely $1-|\text{fiber}|/m$.
The minimum is no larger. Balanced fibers with their complementary
events attain the resulting upper bound. A perfectly informed encoder
can simulate any less informative encoder, so this is an upper bound
at every $\rho$.

For the low-reveal equality let balanced fibers have weights $w_j=n_j/m$.
When $\rho<1$, initialize their label register with probabilities
\begin{equation}\label{eq:integer-initial}
 r_j=\frac{\rho w_j+(1-2\rho)/K}{1-\rho}.
\end{equation}
These sum to one, and are nonnegative precisely throughout the stated
range: the smallest numerator is $\rho q/m+(1-2\rho)/K$.
On a revealed parameter overwrite the label by its fiber; on erasure
do nothing. No revelation flag is retained. For a parameter in fiber
$i$, the complementary-event score is
\begin{align*}
 g_i&=\rho(1-w_i)+(1-\rho)\left(r_i-\sum_jr_jw_j\right)\\
    &=\rho\left(1-\sum_jw_j^2\right)=\rho B_{m,K}.
\end{align*}
This matches the upper bound uniformly in the fixed parameter. If
$\rho=1$ and $s=0$, the deterministic revealed label already attains
that value, so the singular denominator in \eqref{eq:integer-initial}
is not used. Finally, with unequal balanced fibers the weighted average
of their sizes is strictly smaller than the largest size; equivalently
$B_{m,K}>1-\lceil m/K\rceil/m$. This proves the last assertion.
\end{proof}

\begin{theorem}[The complete two-label revelation law]\label{thm:two-label-erasure}
For every $m\ge2$, every $N\ge0$, and $0<\alpha<1$, set
$w=\lfloor m/2\rfloor/m$ and $b=\lceil m/2\rceil/m$. Then
\begin{equation}\label{eq:two-label-erasure}
 V^{\rm er}_{N,m}(2)=\min\{2wb\rho,\;w\},
                 \qquad \rho=1-(1-\alpha)^N.
\end{equation}
The same formula holds at the limiting endpoints $\rho=0,1$.
For odd $m$ it changes regime at $\rho=m/(m+1)$; for even $m$ it
reduces to $\rho/2$. The attaining machine uses two training labels,
a frozen two-event alphabet, and atomic peak at most four.
\end{theorem}
\begin{proof}
Proposition~\ref{prop:integer-erasure} gives the upper bound and its
attainment for $\rho\le1/(2b)$. Above this threshold, partition into
fibers of weights $w,b$ and initialize deterministically in the larger
fiber's label. Continue to overwrite on a reveal and leave the state
unchanged on an erasure. A parameter in the larger fiber has score
$w$; a parameter in the smaller fiber has score $(2\rho-1)b$.
The latter is at least $w$ exactly when $\rho\ge1/(2b)$.
This attains the second upper bound. For even $m$ the threshold is one,
so the low-reveal construction suffices. The endpoint $\rho=0$ may
instead freeze a constant event. Candidate validation retains the
selected label and its indicator, with at most four states; target
validation produces the difference. Initialization is atomic
randomization when its probabilities are not dyadic.
\end{proof}

For example, $m=3,K=2$ gives
$V^{\rm er}_{N,3}(2)=\min\{4\rho/9,1/3\}$, rather than $4\rho/9$
for all $\rho$. The balanced-integer expression is therefore not the
unqualified replacement of the divisible minimax formula. For
$K\ge3$ outside the low-reveal range, \eqref{eq:integer-bounds} remains
a bound rather than an asserted equality. The compatible dynamic
formulation of Section~\ref{sec:controlled-memory} retains the exact
optimization in that regime. These results strengthen the revelation
example without identifying its decision-cut coordinate with the
collision model's bit-level peak.
